\section{Star Differeth from Star}

Heaven is unequal. That is not a pious speculation or a medieval hierarchy imported into the Gospel; it is defined doctrine, in the same sentence that defines the vision itself, and it is one of the points on which the modern devotional picture has drifted furthest from the received teaching without anyone announcing a change.

\subsection{The definition, and the texts behind it}

The Council of Florence in 1439 defined that the purified souls are received at once into heaven and see clearly God himself, three and one, as he is---\latin{pro meritorum tamen diversitate alium alio perfectius}, ``yet, in accordance with the diversity of merits, one more perfectly than another.''\endnote{Council of Florence, \emph{Laetentur caeli}, 6 July 1439, DS 1305; Latin read 2026-07-26 at the registered patristica.net presentation of the Latin \emph{Enchiridion symbolorum}. The English is this article's labelled working gloss.} The Council of Trent, in the decree on justification of 1547, anathematised anyone who denies that a justified person truly merits, by good works done through God's grace and Christ's merit, ``an increase of grace, eternal life and---if he dies in grace---the attainment of that eternal life, and even an increase of glory.''\endnote{Council of Trent, session VI (13 January 1547), canon 32, DS 1582 (older numbering 842): \latin{\ldots non vere mereri augmentum gratiae, vitam aeternam et ipsius vitae aeternae (si tamen in gratia decesserit) consecutionem, atque etiam gloriae augmentum: an.\ s.} Latin read 2026-07-26 at the registered presentation; English gloss labelled. The canon's structure should be noticed: it condemns the denial that the works are \emph{both} God's gifts \emph{and} the justified person's merits. The doctrine is not that some earn what others are given.} \latin{Gloriae augmentum}: an increase of glory, truly merited. The Second Vatican Council states the same fact descriptively of the Church in its three states---the pilgrims, those being purified, and those in glory---saying that ``all in various ways and degrees are in communion in the same charity of God and neighbor.''\endnote{Second Vatican Council, dogmatic constitution \emph{Lumen gentium} (21 November 1964), n.~49, official English text at vatican.va, read 2026-07-26. The same paragraph quotes the Florence definition---the blessed are ``in glory beholding `clearly God Himself triune and one, as He is'\,''---and cites \emph{Benedictus Deus} at n.~49 as well.}

Behind the definitions stands one text and, in the tradition's actual argumentation, mainly one. Paul, in the middle of the resurrection chapter: ``One is the glory of the sun, another the glory of the moon, and another the glory of the stars. For star differeth from star in glory. So also is the resurrection of the dead.''\endnote{1 Cor 15:41--42, tracked verse text of the Douay--Rheims, read 2026-07-26.} It is that verse, not the more famous ``In my Father's house there are many mansions,'' that Aquinas puts in the \latin{sed contra} when he asks whether one of the blessed sees God more perfectly than another.\endnote{ST I, q.~12, a.~6, \latin{sed contra}, read 2026-07-26 at New Advent: ``Eternal life consists in the vision of God, according to John 17:3\ldots\ Therefore if all saw the essence of God equally in eternal life, all would be equal; the contrary to which is declared by the Apostle: `Star differs from star in glory' (1 Corinthians 15:41).'' The many-mansions verse (Jn 14:2) has been read as teaching diverse degrees since antiquity, but the sense of \latin{mansiones} there is genuinely disputed---dwelling-places, or stages, or simply room enough---and this article does not rest the claim on it.}

\subsection{The mechanism, and what it excludes}

Article six of the twelfth question is short and its argument is worth following exactly, because what it rules out is more interesting than what it rules in.

The difference cannot be on the side of the object: God is one, and the blessed do not see him through a likeness that could be sharper in one mind than another. ``The same object will be presented to all---viz.\ the essence of God.''\endnote{ST I, q.~12, a.~6, reply to objection 3, read 2026-07-26.} So no one in heaven has a better God than anyone else, and nothing is withheld from the lesser. That single point disposes of most of what makes the doctrine of degrees sound repellent.

Nor can the difference be a difference of natural intellectual power. Aquinas raises the objection himself and treats it as decisive: if the sharper mind saw more, then rank in glory would follow rank in nature, ``and this is incongruous; since equality with angels is promised to men as their beatitude.''\endnote{ST I, q.~12, a.~6, objection 3, read 2026-07-26. The reply concedes the principle and locates the difference in ``the glorified faculty'' rather than the natural one.} Heaven is not graded by intelligence, and it is not graded by learning, office, or anything else a natural endowment could measure.

What remains is the light of glory, and the light is measured by charity:

\begin{studyframe}
``Hence the intellect which has more of the light of glory will see God the more perfectly; and he will have a fuller participation of the light of glory who has more charity; because where there is the greater charity, there is the more desire; and desire in a certain degree makes the one desiring apt and prepared to receive the object desired. Hence he who possesses the more charity, will see God the more perfectly, and will be the more beatified.''\endnote{ST I, q.~12, a.~6, corpus, read 2026-07-26 at New Advent; Latin at Corpus Thomisticum: \latin{Unde intellectus plus participans de lumine gloriae, perfectius Deum videbit. Plus autem participabit de lumine gloriae, qui plus habet de caritate, quia ubi est maior caritas, ibi est maius desiderium; et desiderium quodammodo facit desiderantem aptum et paratum ad susceptionem desiderati.}}
\end{studyframe}

The mechanism, then, is not payment but capacity. Desire widens what desires. A life of charity does not accumulate credit to be exchanged at the end for a better seat; it enlarges the person, and the enlargement is what shows in glory. This is why the doctrine of degrees and the doctrine of grace do not collide: the charity that measures the capacity is itself infused, so the whole scale is a scale of gifts received, and Trent's canon is careful to say in one breath that the works are God's gifts and the man's merits.

\subsection{How equality and difference are held together}

Three things hold the doctrine together, and all three are in the sources rather than in the apologetics.

\emph{The object is common.} Everyone in the city has God, and has him immediately, and has all of him---no one is looking at a smaller God. What differs is the seeing, not the seen. Aquinas is explicit that the vision is not by a likeness at all, so there is no sense in which the lesser blessed have a blurrier picture of the same thing; they have less capacity for an object that is wholly given to each.

\emph{Everyone is full.} The tradition's way of putting this comes from Augustine, in the chapter on the eternal felicity of the city of God, and it is worth quoting at length because it is the classical answer:

\begin{studyframe}
``But who can conceive, not to say describe, what degrees of honour and glory shall be awarded to the various degrees of merit? Yet it cannot be doubted that there shall be degrees. And in that blessed city there shall be this great blessing, that no inferior shall envy any superior, as now the archangels are not envied by the angels, because no one will wish to be what he has not received, though bound in strictest concord with him who has received; as in the body the finger does not seek to be the eye, though both members are harmoniously included in the complete structure of the body. And thus, along with his gift, greater or less, each shall receive this further gift of contentment to desire no more than he has.''\endnote{Augustine, \emph{De civitate Dei} XXII.30, trans.\ Marcus Dods (1871), public domain; read 2026-07-26 in the source library's tracked transcription of the 1871 printing at the registered passage for XXII.30. Note that Augustine begins with ``who can conceive, not to say describe'' and only then says the degrees cannot be doubted: the fact is asserted, the content disclaimed, in a single breath.}
\end{studyframe}

Notice what Augustine does and does not claim. He asserts the fact of degrees as beyond doubt and declares their content inconceivable in the same sentence. And he adds a second gift on top of the first: contentment, given with the measure, so that the measure cannot be experienced as a deprivation. Whether that is a satisfying answer to a modern reader's sense of justice is a fair question. It is at any rate the answer, and it is not the answer of a man who thinks heaven is a prize table.

\emph{The scale is charity, which is nobody's achievement.} If the ranking ran on cleverness, or courage, or public service, it would be a ranking of things people can compare. It runs on the love of God, which is infused, which is hidden, and which no one can measure in himself, let alone in anyone else. The doctrine of degrees, correctly stated, ends every conversation about who is greater rather than starting one.

\subsection{The Eastern way of holding the same thing}

Gregory of Nyssa's account offers a different reconciliation, and it is worth setting beside the Latin one rather than against it. On his account the soul is a vessel enlarged by what fills it, so that ``he never stops enlarging''; the capacity that measures beatitude is not a fixed allotment but a growth without limit.\endnote{Gregory of Nyssa, \emph{On the Soul and the Resurrection}, NPNF 2nd series, vol.~5 (Moore), public domain, read 2026-07-26 at New Advent; quoted more fully in the section on the Eastern answer above.} If that is right, then the difference between the greater and the lesser blessed is not a permanent rank but a place in a motion that never stops---which dissolves the offence of the doctrine at the cost of raising a question the Latin tradition would want answered, namely how an unending increase squares with the rest that \emph{Benedictus Deus} says the blessed already have.

This article does not adjudicate that, and notes only that the two accounts are not obviously incompatible: an unchanging vision of an infinite object and an unending enlargement of the seer are different pictures of the same situation, and nothing defined by the Catholic Church rules out the second. What is defined is the inequality and its measure, not its permanence in every respect.

\subsection{What is speculation here}

A good deal, and it is easy to spot once the criterion is applied.

The doctrine of the aureoles---the additional crowns of virgins, martyrs, and doctors, with the question of which is greatest---is school theology built on the essential-versus-accidental-reward distinction. The distinction is reasonable and the elaboration is not defined.\endnote{ST Supplement, q.~96, aa.~1, 5--7, 12, headings and a.~1 corpus read 2026-07-26 at New Advent; see the section on the risen body above for the provenance of the Supplement.} Any assignment of persons to degrees is speculation of a worse kind, since it requires knowledge of interior charity that no one has. The idea that degrees can be inferred from visible achievement is contradicted by the mechanism itself. And the confident modern reflex in the other direction---that heaven must be equal because inequality would be unfair---is not a correction of a medieval excess but a rejection of a conciliar definition, usually made by people who do not know they are making it.

The honest position is narrow and can be stated in three sentences. The blessed see God unequally, in proportion to charity, and this is defined. The inequality is a difference of capacity for an object wholly given to each, and this is the common and convincing explanation, not a definition. What any particular degree consists in, how it is experienced, whether it is fixed or grows, and where anyone stands on it, is unknown, and the greatest Western authority on the subject said so in the sentence in which he affirmed the doctrine.
